\clearpage
\section{总结}

\subsection{对学习方法的直观认识}

本次实验让我们直观体验到了基于学习的单目深度预测方法已经能够达到较好的效果。仅输入一张普通的 RGB 图像，Depth Anything V2 就能预测出场景中较清晰的前后关系，区分墙面、地面和家具，并保留部分细小结构。经过深度对齐和点云生成后，我们还能够从不同角度观察重建出的三维场景。这让我们认识到，模型从数据中学到的经验能够为深度恢复提供有效线索，也让课上的三维重建过程变得更加直观。

\subsection{对补充深度信息的理解与后续思考}

实验也加深了我们对课上所学内容的理解：\textbf{从二维图像恢复三维空间时，像素位置只能确定一条射线，缺少的深度信息需要通过其他途径补充。}本次实验中，模型先给出相对深度结构，再利用真实深度标注校正尺度与偏移。真实标注提供了较可靠的尺度约束，使预测结果能够以米为单位表示，并生成与参考场景较接近的点云。Depth Anything V2 (Metric Depth) 在对齐后误差明显下降，也让我们看到了这类补充信息对重建效果的帮助。

因此，我们认为，后续要继续提高三维恢复效果，除了改进模型，也可以寻找更多可靠的方式来\textbf{补充深度约束}。例如，让机器人主动移动，从不同位置拍摄同一场景，形成类似双目的多视角观测；结合相机位姿和已知的移动距离，利用视差估计深度。这样就能把模型学到的场景经验与实际观测得到的几何信息结合起来，尝试在没有真实深度标注的情况下，获得更准确、具有真实尺度的三维重建。
